% Blocos LaTeX para integrar no Capitulo 3.
% Sugestao de uso:
% 1) Inserir a nova subseccao "Metodo Multimodal com PaddleOCR-VL" na seccao
%    "Implementacao dos Metodos de Extracao", apos "Metodo OCR + LLM".
% 2) Substituir a tabela "Comparacao sintetica..." pela versao atualizada abaixo.
% 3) Inserir os pequenos paragrafos de atualizacao nas subseccoes indicadas.

% -------------------------------------------------------------------------
% Acrescentar em "Enquadramento Metodologico", no paragrafo que enumera os
% pipelines alternativos, apos a referencia ao metodo OCR + LLM.
% -------------------------------------------------------------------------

Numa fase posterior, foi tambem integrado um quinto metodo, baseado em
PaddleOCR-VL, com o objetivo de avaliar uma abordagem multimodal de parsing
documental. Ao contrario dos pipelines assentes apenas em OCR textual ou em
detecao visual supervisionada, este metodo utiliza um modelo visao-linguagem
especializado em documentos, capaz de reconhecer simultaneamente texto, layout,
tabelas, formulas e outros elementos visuais. Esta integracao permitiu testar
uma familia de abordagem mais proxima dos modelos recentes de document parsing,
mantendo, contudo, a mesma estrutura de avaliacao e comparacao usada nos
restantes metodos.

% -------------------------------------------------------------------------
% Acrescentar em "Arquitetura Geral da Solucao", apos a lista de modulos ou
% depois do paragrafo sobre backend/frontend.
% -------------------------------------------------------------------------

A integracao do metodo PaddleOCR-VL introduziu ainda uma variante arquitetural
semi-automatizada, na qual a inferencia pesada do modelo multimodal e executada
em Google Colab, recorrendo a GPU, enquanto o backend local mantem a
responsabilidade de importar, normalizar, avaliar e apresentar os resultados.
Nesta configuracao, o PDF continua a ser o documento de referencia submetido no
frontend, mas o utilizador associa tambem o ficheiro \texttt{run\_summary.json}
produzido no Colab. Este ficheiro contem os resultados brutos do PaddleOCR-VL
por pagina, incluindo blocos detetados, coordenadas, etiquetas de layout,
conteudo textual e tabelas em HTML. O backend converte posteriormente esta saida
para o esquema comum da aplicacao, garantindo que o metodo multimodal pode ser
comparado com os restantes pipelines.

% -------------------------------------------------------------------------
% Inserir em "Implementacao dos Metodos de Extracao", depois de OCR + LLM.
% -------------------------------------------------------------------------

\subsection{Metodo Multimodal com PaddleOCR-VL}

Para alem dos metodos baseados em OCR classico, extracao direta de tabelas,
detecao visual supervisionada e interpretacao por LLM, foi integrado um metodo
experimental baseado em PaddleOCR-VL. Este metodo foi introduzido com o objetivo
de avaliar uma abordagem multimodal especificamente concebida para parsing
documental, isto e, para leitura conjunta da imagem da pagina, da sua estrutura
visual e dos elementos textuais ou tabulares nela presentes.

O PaddleOCR-VL e um modelo visao-linguagem compacto orientado para documentos.
De acordo com a documentacao e o relatorio tecnico do modelo, o seu nucleo e o
PaddleOCR-VL-0.9B, que combina um codificador visual de resolucao dinamica do
tipo NaViT com o modelo de linguagem ERNIE-4.5-0.3B
\cite{cui2025paddleocrvl,paddleocrvl_docs}. A versao PaddleOCR-VL-1.5 mantem a
logica de um VLM compacto de 0.9 mil milhoes de parametros, alargando a robustez
em cenarios documentais reais e tarefas como parsing, text spotting e
reconhecimento de selos \cite{cui2026paddleocrvl15}. Esta arquitetura e
particularmente relevante para documentos semi-estruturados porque permite
processar simultaneamente informacao visual e textual, em vez de tratar a pagina
apenas como uma sequencia plana de caracteres.

Em termos funcionais, o pipeline do PaddleOCR-VL combina uma etapa de analise de
layout com uma etapa de reconhecimento multimodal. A analise de layout identifica
regioes da pagina, como cabecalhos, texto, tabelas, formulas, imagens, rodapes
ou outros blocos visuais. Em seguida, o modelo visao-linguagem interpreta cada
regiao e produz conteudo estruturado, incluindo texto normalizado, tabelas em
HTML e formulas em representacao matematica quando aplicavel. A saida conserva
tambem metadados importantes, como a pagina de origem, a classe do bloco e a
caixa delimitadora, permitindo rastrear a origem visual de cada elemento
extraido.

No prototipo desenvolvido, este metodo foi executado em Google Colab devido aos
requisitos computacionais associados ao modelo. Os PDFs sao inicialmente
convertidos em imagens, organizadas por documento, e processados com o comando
\texttt{paddleocr doc\_parser}. O fluxo implementado permite executar todas as
pastas presentes no diretorio de entrada, produzindo um \texttt{run\_summary.json}
por documento. Este resumo contem, para cada pagina, o estado da execucao, os
logs do processo e a representacao dos blocos extraidos pelo modelo. Para reduzir
o custo temporal associado ao carregamento repetido dos pesos, foi ainda
preparado um modo de execucao por documento, que evita recarregar o modelo
pagina a pagina sempre que o ambiente o permite.

A integracao com a aplicacao local foi desenhada de forma a preservar a
arquitetura comparativa ja existente. No frontend, quando o metodo
\texttt{paddleocr\_vl} e selecionado, o utilizador submete o PDF e associa o
ficheiro \texttt{run\_summary.json} produzido no Colab. O backend guarda ambos os
ficheiros, copia o resumo para a estrutura esperada em
\texttt{data/paddleocr\_vl/<documento>/run\_summary.json} e executa o script
\texttt{15\_import\_paddleocr\_vl.py}. Este script extrai a lista de blocos
produzida pelo PaddleOCR-VL, remove elementos de log, interpreta tabelas HTML,
normaliza labels e converte a informacao para os mesmos ficheiros usados pelos
outros metodos: \texttt{parsed}, \texttt{tables}, \texttt{sections} e
\texttt{metrics}.

Uma parte importante desta integracao consistiu em evitar regras estaticas
dependentes do nome do documento. Em vez disso, foram implementadas regras
genericas baseadas em labels documentais e estrutura tabular. Por exemplo, linhas
com etiquetas como \texttt{Nome}, \texttt{Morada}, \texttt{Designacao},
\texttt{Marca}, \texttt{Modelo}, \texttt{N. Serie}, \texttt{Local},
\texttt{Temperatura} ou \texttt{Humidade} sao interpretadas de forma transversal,
independentemente do certificado concreto. Do mesmo modo, tabelas com cabecalhos
associados a \texttt{Padrao}, \texttt{CATIM} e \texttt{Rastreabilidade} sao
tratadas como tabelas de equipamento padrao, enquanto tabelas com colunas de
forca, erro, fator de expansao e incerteza sao convertidas em tabelas de
resultados de calibracao de forca. As restantes tabelas reconhecidas pelo modelo
sao preservadas como tabelas genericas detetadas, permitindo auditoria mesmo
quando ainda nao existe uma interpretacao semantica especializada.

Este metodo revelou-se uma boa escolha para o contexto dos certificados de
calibracao por varias razoes. Em primeiro lugar, os documentos analisados
apresentam forte heterogeneidade visual, com variacoes de layout, tabelas densas,
formulas, rodapes extensos e informacao administrativa distribuida por blocos. Um
modelo multimodal especializado em documentos e mais adequado a esta realidade do
que um OCR de pagina completa, pois considera explicitamente a organizacao visual
da pagina. Em segundo lugar, os certificados podem conter portugues, espanhol,
ingles tecnico e simbologia metrologica; o suporte multilingue e a capacidade de
reconhecer elementos complexos tornam o PaddleOCR-VL particularmente relevante
para este cenario. Em terceiro lugar, por ser um modelo relativamente compacto,
permite experimentacao em ambientes acessiveis, como Google Colab, sem depender
necessariamente de APIs comerciais.

Apesar destas vantagens, a abordagem tambem apresenta limitacoes. O tempo de
execucao e superior ao de pipelines OCR mais simples, sobretudo quando o modelo
e carregado repetidamente. A instalacao das dependencias e mais pesada e a
execucao em Windows nao e tao direta como em Linux ou ambientes containerizados.
Adicionalmente, nem todas as regioes relevantes sao sempre devolvidas como
tabelas estruturadas; em alguns certificados, zonas de resultados podem ser
classificadas como grafico ou imagem, exigindo uma etapa futura de fallback ou
recorte adicional. Assim, o PaddleOCR-VL deve ser entendido como um metodo
multimodal promissor e complementar, nao como substituto imediato de toda a
arquitetura experimental.

Metodologicamente, a sua inclusao e relevante porque permite comparar tres
niveis distintos de compreensao documental: OCR textual de pagina completa,
segmentacao visual supervisionada com YOLO e parsing multimodal com um VLM
especializado. Esta comparacao e particularmente util para perceber ate que
ponto modelos gerais de document parsing conseguem reduzir a necessidade de
templates e anotacao manual, mantendo a rastreabilidade e a possibilidade de
normalizacao para o esquema final do SIMME.

% -------------------------------------------------------------------------
% Acrescentar em "Metricas e Estrategia de Avaliacao", apos a explicacao do
% TableExtractionScore.
% -------------------------------------------------------------------------

Com a integracao do PaddleOCR-VL, tornou-se tambem necessario distinguir entre
duas nocoes de desempenho tabular. A primeira corresponde ao
\textit{TableExtractionScore}, calculado apenas quando existe um tipo de
instrumento com tabelas esperadas conhecidas. A segunda corresponde a deteccao
generica de grupos tabulares, usada quando o metodo devolve tabelas estruturadas
mas ainda nao existe uma correspondencia formal com uma tabela esperada do
esquema. Para esse efeito, foram registados adicionalmente os contadores
\texttt{row\_counts} e \texttt{detected\_tables}. Esta distincao evita interpretar
um caso \texttt{0/0} como ausencia de tabelas, permitindo reconhecer situacoes em
que o modelo extraiu tabelas uteis, embora ainda nao classificadas como tabela de
resultados do instrumento.

% -------------------------------------------------------------------------
% Substituir a Tabela "Comparacao sintetica..." pela versao abaixo.
% -------------------------------------------------------------------------

\begin{table}[htbp]
\centering
\caption{Comparacao sintetica entre os metodos implementados no prototipo}
\label{tab:metodosimplementados}
\scriptsize
\setlength{\tabcolsep}{3pt}
\begin{tabular}{p{2.3cm} p{2.0cm} p{3.1cm} p{3.0cm} p{3.1cm}}
\hline
\textbf{Metodo} & \textbf{Entrada} & \textbf{Processamento e saida} & \textbf{Vantagens} & \textbf{Limitacoes} \\
\hline
\texttt{paddle\_current} & PDF convertido em imagem & OCR de pagina completa, segmentacao e parsing heuristico para JSON com campos e tabelas & Simplicidade, baixo custo computacional e elevada auditabilidade & Sensivel a ruido visual, degradacao em layouts densos e dificuldade em reconstrucao estrutural \\
\texttt{pdf\_table} & PDF nativo & Extracao direta de tabelas com \textit{pdfplumber}, com JSON de tabelas e campos inferidos & Util em PDFs digitais limpos e tabulares; rapido de executar & Fragil em documentos pouco estruturados, com tabelas implicitas ou PDFs visualmente complexos \\
\texttt{hybrid\_fast} & Imagens das paginas & Detecao de regioes com YOLO seguida de OCR localizado com PaddleOCR & Melhor isolamento semantico das zonas relevantes e maior robustez ao ruido visual & Maior complexidade de implementacao; requer anotacao manual e treino supervisionado \\
\texttt{ocr\_llm} & Imagens das paginas e texto OCR & OCR combinado com interpretacao semantica por LLM para JSON normalizado & Flexibilidade perante layouts heterogeneos e variacao terminologica & Maior variabilidade, custo potencial e necessidade de validacao adicional \\
\texttt{paddleocr\_vl} & Imagens das paginas e \texttt{run\_summary.json} gerado no Colab & Parsing multimodal com PaddleOCR-VL, importacao de blocos, tabelas HTML e metadados visuais para o esquema comum & Reconhecimento conjunto de layout, texto, tabelas e formulas; adequado a documentos heterogeneos e multilingues; reduz dependencia de templates & Execucao mais pesada; depende de GPU/Colab nesta fase; pode devolver algumas regioes como imagem ou grafico em vez de tabela estruturada \\
\hline
\end{tabular}
\end{table}

% -------------------------------------------------------------------------
% Acrescentar em "Tecnologias Utilizadas", na lista de tecnologias.
% -------------------------------------------------------------------------

\begin{itemize}
    \item PaddleOCR-VL, como modelo multimodal de parsing documental usado para
    reconhecer layout, texto, tabelas, formulas e outros elementos visuais;
    \item Google Drive, como mecanismo de transferencia dos ficheiros de entrada
    e saida entre o Colab e a aplicacao local no metodo PaddleOCR-VL.
\end{itemize}

% -------------------------------------------------------------------------
% Acrescentar em "Estrutura e Funcionamento do Backend", apos a explicacao
% sobre registo de metodos e execucao dos scripts.
% -------------------------------------------------------------------------

No caso especifico do metodo \texttt{paddleocr\_vl}, o backend recebeu uma
adaptacao adicional para aceitar, alem do PDF, o ficheiro
\texttt{run\_summary.json} produzido no Google Colab. Este ficheiro e validado no
momento do upload e associado ao mesmo \texttt{file\_id} do PDF. Antes da
execucao do pipeline, o backend copia o resumo para a pasta
\texttt{data/paddleocr\_vl}, preservando a associacao ao nome do documento. Em
seguida, o script de importacao converte a saida multimodal para o formato
interno da aplicacao. Desta forma, a execucao do modelo pesado pode ocorrer fora
do ambiente local, mas o resultado continua a ser integrado na mesma API, nas
mesmas metricas e na mesma interface de visualizacao.

% -------------------------------------------------------------------------
% Acrescentar em "Estrutura e Funcionamento do Frontend".
% -------------------------------------------------------------------------

Para suportar o metodo PaddleOCR-VL, o frontend foi igualmente adaptado. Quando
este metodo e selecionado, o formulario de upload passa a solicitar tambem o
ficheiro \texttt{run\_summary.json} gerado no Colab. A visualizacao dos
resultados foi alargada para mostrar nao apenas os campos extraidos, mas tambem
grupos tabulares genericos detetados pelo modelo. Foi ainda mantido um bloco de
JSON bruto, com opcao de copia, que permite auditar a saida completa e verificar
a correspondencia entre blocos originais, tabelas interpretadas e metricas.

% -------------------------------------------------------------------------
% Acrescentar em "Observacoes Experimentais Preliminares", antes do paragrafo
% final sobre OCR + LLM ou depois dele.
% -------------------------------------------------------------------------

O metodo \texttt{paddleocr\_vl} apresentou comportamento promissor na leitura de
certificados com estrutura visual complexa. Em particular, mostrou capacidade
para devolver blocos semanticos com etiquetas de layout, reconhecer tabelas em
HTML e preservar informacao de coordenadas, o que facilita a auditoria posterior.
Em certificados da familia de calibracao de forca, permitiu recuperar campos
administrativos, condicoes de trabalho, equipamento padrao utilizado e, quando a
regiao foi reconhecida como tabela, resultados tabulares de calibracao. No
entanto, observou-se tambem que algumas regioes de resultados podem ser
classificadas como \textit{chart} ou imagem, nao sendo imediatamente convertidas
em linhas tabulares. Esta observacao reforca a necessidade de combinar a saida
multimodal com camadas adicionais de normalizacao e fallback, em vez de depender
exclusivamente da saida direta do modelo.

% -------------------------------------------------------------------------
% Substituir/expandir a Sintese do Capitulo com o paragrafo seguinte.
% -------------------------------------------------------------------------

A integracao do PaddleOCR-VL veio acrescentar uma dimensao multimodal a esta
comparacao. Enquanto o baseline e o metodo \textit{pdfplumber} representam
abordagens mais tradicionais, e o metodo YOLO + PaddleOCR representa uma solucao
hibrida supervisionada, o PaddleOCR-VL permite avaliar um modelo especializado
em parsing documental com capacidade nativa para interpretar layout e conteudo.
Os resultados preliminares sugerem que esta abordagem e particularmente
interessante para documentos heterogeneos e ricos em estrutura visual, embora
ainda exija pos-processamento para garantir a normalizacao completa dos campos e
tabelas necessarios ao SIMME.

